% vim: set spell:
% vim: set spell spelllang=en:
\documentclass[a4paper,10pt]{article}
%\usepackage[mathjax]{lwarp}
%\usepackage{lwarp}
\usepackage[utf8]{inputenc}
\usepackage[T1]{fontenc}
\usepackage{amsfonts}
\usepackage{amsmath}
\usepackage{pifont}
\usepackage{tikz}
\usetikzlibrary{calc}
\usepackage{hyperref}
\usepackage{stmaryrd}
\usepackage[left=2cm,right=2cm]{geometry}

\title{Tutorials : Concepts of Logic and Programming}
\author{Nicolas González}
\date{September 9th 2020}

\def\centralruler{\begin{center}\rule{5cm}{.1pt}\end{center}}

\newcommand{\R}{\mathbb{R}}
\newcommand{\Com}{\mathbb{C}}
\newcommand{\K}{\mathbb{K}}
\newcommand{\Z}{\mathbb{Z}}
\newcommand{\N}{\mathbb{N}}
\newcommand{\Lin}{\mathcal{L}}
\newcommand{\Mnn}{\mathcal{M}_n(\R)}
\newcommand{\ProScal}[2]{\langle #1,#2 \rangle}
\newcommand{\Norm}[1]{\left|\left| #1 \right|\right|}
\newcommand{\grad}{\overrightarrow{\nabla}}
\newcommand{\Proba}{\mathbb{P}}
\newcommand{\vv}[1]{\overrightarrow{\textbf{#1}}}

\DeclareMathOperator{\Vect}{Vect}
\DeclareMathOperator{\Imag}{\Im m}
\DeclareMathOperator{\Mat}{Mat}
\DeclareMathOperator{\rk}{rk}
\DeclareMathOperator{\asinh}{asinh}


\renewcommand{\theequation}{\thesubsection--\arabic{equation}}


\newcommand{\Tstrut}{\rule{0pt}{2.6ex}}
\newcommand{\Bstrut}{\rule[-0.9ex]{0pt}{0pt}}
\newcommand{\TBstrut}{\Tstrut\Bstrut}
%\usepackage[toc]{glossaries}
%\makeglossaries
%\loadglsentries{glossary}

\begin{document}

\maketitle
%\tableofcontents

\newpage

\section{Exercise 1}
\begin{enumerate}
	\item Convertir les nombres binaires suivants en décimal, octal, hexadecimal
	\begin{itemize}
		\item 0011001
		\item 100001
		\item 0101101
		\item 11010
		\item 1111111
	\end{itemize}
	\item Représentez au format binaire les nombres en base 10 suivants : $(12)_{10}$ $(1025)_{10}$.
\end{enumerate}
\centralruler

\subsection{Question 1}

Let's begin by converting the first number to each of these bases to show how that's done. I'll just skip and give the results at the end.

Notice that the binary notation works with powers of 2. In fact, 0011001 means
$$2^0+2^3+2^4=1+8+16=25$$
That's a small enough number that I can tell you it's also equal to
$$25=16 + 9 = 1*16^1 + 9 * 16^0 = (19)_{16}$$
And again with base $8$
$$25=24+1 = 3*8+1 = 3*8^1+1*8^0=(31)_{8}$$

And there we go. I've summed up all the numbers in a table but the method works every time.

\begin{center}
	\begin{tabular}{|c|c|c|c|}
		\hline
		Base 2&Base 10&Base 16&Base 8\\\hline
		0011001&25&19&31\\\hline
		100001&33&21&41\\\hline
		0101101&45&2D&55\\\hline
		11010&26&1A&32\\\hline
		1111111&127&7F&177\\\hline
	\end{tabular}
\end{center}

If you're faced with larger numbers you might want to use the method based on long (read: euclidian) divisions.

\subsection{Question 2}

Here, knowing the powers of 2 will help us tremendously. For example, knowing that $2^{10}=1024$, so $1025=1024+1$, yields immediately
\[
(1025)_{10} = 0b10000000001
\]

Similarly, $12=8+4 = 2^3+2^2$ so $(12)_{10}=0b1100$.

\section{Exercice 2}
\begin{enumerate}
	\item Exprimez $(+63)_{10}$ et $(-63)_{10}$ en complément à 2 pour un codage sur un octet.
	\item Réalisez la somme en complément à 2 de $(30)_{10}+(-8)_{10}$.
\end{enumerate}
\centralruler

\subsection{Question 1}

Reprensenting positive integers is trivial. For example :
$$63 = 64 - 1 = 2^6 - 1 = 00111111$$
So we can code everything with 8 bits, including the inverse. Indeed, with a complement to $2$ notation, the range of integers representable is
$$[2^{n-1};2^{n-1}-1]$$
Here $n=8$ since our largest is $2^7-1$. So
$$11000000 = -2^{6} = -128 + (64) = -64$$
Which means that
$$11000001 = -2^{6}+1 = -63$$

\subsection{Question 2}

Note that all of our integers here belong in $[-2^5;2^5-1]=[-32;31]$, so let's encode
$$30 = 32 - 2 = 011110$$
$$-8 - (-32) = 24 = 16 + 8 = 011000$$
Meaning that
$$-8 = (111000)_{2}$$
So we can work out the sum of these two numbers in binary
$$(111000)_2+(011110)_2=(010110)_2 = 22$$

\section{Exercice 3}

A l'aide des circuits présentés jusqu'ici, réaliser un circuit qui, à partir de deux opérandes $4$ bits X, Y et de commandes en nombres suffisants, délivre en résultat soit $X+Y$, soit $2X+Y$, soit $X+2Y$
Temporary because the file is on the windows laptop

Exercises 3 through 5 included are provided as logisim files.

\section{Exercise 6}

On veut réaliser la fonction suivante : $S=1$ si $N=0$ ou $3$ ou $5$ ou $7$, $N$ étant un nombre codé sur 3 bits (C, B, A).
\begin{itemize}
	\item Écrire la table de vérité de $S$
	\item Réaliser la fonction avec le multiplexeur adéquat
	\item De façon générale, combien existe-t-il de fonctions différentes de $M$ variables ? En déduit l'intérêt général du multiplexeur.
\end{itemize}

\centralruler

\subsection{Question 1}

Our output $S$ will only be true when $N=0$, or $N=3$, or $N=5$, or $N=7$, the binary patterns for those are (when using C, B, A in decreasing weight) :
\begin{center}
	\begin{tabular}{|c|c|c|c|}
		\hline
		Decimal&C&B&A\\\hline
		0&0&0&0\\\hline
		3&0&1&1\\\hline
		5&1&0&1\\\hline
		7&1&1&1\\\hline
	\end{tabular}
\end{center}

The truth table becomes
\begin{center}
	\begin{tabular}{|c|c|c|c|}
		\hline
		C&B&A&Output\\\hline
		0&0&0&1\\\hline
		0&0&1&0\\\hline
		0&1&0&0\\\hline
		0&1&1&1\\\hline
		1&0&0&0\\\hline
		1&0&1&1\\\hline
		1&1&0&0\\\hline
		1&1&1&1\\\hline
	\end{tabular}
\end{center}

The truth function for that should be something like
\[
	S=\overline{A}~\overline{B}~\overline{C}+A~(B+C)
\]

\subsection{Question 2}

This question is done in Logisim.

\subsection{Question 3}

Generally speaking, from a purely combinatronics point of view, for $M$ variables, there can exist exactly $2^M$ functions. When those function have no meaningful sense, or one that is so hard that it would require an impossible amount of gates, we can use a multiplexer to simply create a mapping between an input number and a binary output value.

\section{Exercise 7}

On se propose de réaliser un "contrôleur de vumètre" : c'est un dispositif qui permet de visualiser une grandeur numérique (un volume sonore par exemple) au moyen d'une barrette lumineuse qui s'étire plus ou moins en fonction de la grandeur.

\begin{enumerate}
	\item Réaliser, par la méthode booléenne, le contrôleur de vumètre illustré ci-dessus. La grandeur à visualiser est représentée en binaire sur 3 bits ; la barrette est constituée de 7 segments qui s'allument un par 1. Pour la grandeur 0, la barrette doit être entièrement éteinte. Réaliser ce même circuit en utilisant judicieusement un décodeur (et des portes).
	\item Réaliser un circuit similaire pour une entrée 4 bits, avec une barrette de 15 segments.
	\item Réaliser un circuit similaire pour une grandeur relative (positive ou négative) représentée en complément sur 4 bits, la barrette étant divisée en les 7 segments du haut pour les grandeurs positives et les 8 segments du bas pour les grandeurs négatives.
\end{enumerate}

\centralruler

\subsection{Question 1}
It's easier to consider the individual segments out. For example, the first one is on if any of the input bits is true. So
\[
	s1=v0+v1+v2
\]
For the second one, we should at least have the second (or third) bit on, so
\[
	s2=v1+v2
\]
For the third segment, either both $v0$ and $v1$ are on, or $v2$ is.
\[
	s3=v0~v1+v2
\]
For the fourth segment, and all those coming after it, $v2$ is necessarily on.
\[
	s4=v2
\]
For the fifth segment, $v2$ is on with at least $v0$ or $v1$ on.
\[
	s5=v2~(v1+v0)
\]
For the sixth segment, $v2$ is on but so must be $v1$.
\[
	s6=v2~v1
\]
The seventh segment is only on when all of the input bits are on.
\[
	s7=v0~v1~v2
\]

And this is all good, but very redundant. See, if a segment is on, all of those before it should be on, so, if we use a decoder, we can propagate its output using OR gates and get the same result without the hassle of a billion and gates.

I've done that in the Logisim file for this exercise.

\subsection{Question 2}

Done in Logisim.

\subsection{Question 3}

Also done in Logisim. (same file)
\end{document}
